\documentclass[aps,prd,twocolumn,showpacs,superscriptaddress,nofootinbib]{revtex4-2}

\usepackage{amsmath,amssymb}
\usepackage{graphicx}
\usepackage{booktabs}
\usepackage{siunitx}
\usepackage{hyperref}

\begin{document}

\title{Spectral Predictions for Analog Hawking Radiation:\\
Platform-Specific Tables from SK-EFT}

\author{SK-EFT Hawking Research Program}

\date{\today}

\begin{abstract}
We present platform-specific spectral predictions for analog Hawking
radiation in three BEC experiments: Steinhauer $^{87}$Rb (Technion),
Heidelberg $^{39}$K, and Trento $^{23}$Na spin-sonic. The predictions
derive from the Schwinger-Keldysh dissipative effective field theory
framework (Papers~1--4) and include: (i)~complete spectral occupation
$n(\omega)$ at representative frequencies, (ii)~correction decomposition
into dispersive and dissipative contributions, (iii)~noise floor
thresholds from the fluctuation-dissipation relation, (iv)~detector
sensitivity requirements for three measurement goals, and
(v)~a kappa-scaling test prediction for Heidelberg~$^{39}$K.
This document is designed for experimentalists planning measurements.
\end{abstract}

\maketitle

\section{Platform Parameters}

\begin{table}[h]
\caption{BEC platform parameters in natural units ($c_s = 1$, $\kappa = 1$).}
\label{tab:params}
\begin{ruledtabular}
\begin{tabular}{lccc}
Parameter & Steinhauer & Heidelberg & Trento \\
\hline
Atom & $^{87}$Rb & $^{39}$K & $^{23}$Na \\
$D$ & 0.013 & 0.012 & 0.014 \\
$\gamma_\text{dim}$ & $2.4\times10^{-5}$ & $1.6\times10^{-3}$ & $1.4\times10^{-5}$ \\
$T_H$ & $\kappa/2\pi$ & $\kappa/2\pi$ & $\kappa/2\pi$ \\
$\omega_\text{max}/T_H$ & 108 & 118 & 108 \\
\end{tabular}
\end{ruledtabular}
\end{table}

\section{Spectral Predictions}

Table~\ref{tab:spectral} lists the predicted occupation number $n(\omega)$
at five representative frequencies for each platform, along with the
decomposition into dispersive ($\delta_\text{disp}$) and dissipative
($\delta_\text{diss}$) corrections.

The total occupation is:
\begin{equation}
n(\omega) = n_\text{Hawking}(\omega) + n_\text{noise}(\omega)
\end{equation}
where $n_\text{Hawking} = |\beta|^2 / (1 - \delta_k)$ includes the
modified unitarity correction and $n_\text{noise} = \delta_k/2$ is the
FDR-mandated noise floor.

\begin{table*}[t]
\caption{Spectral predictions at representative frequencies. Values computed
from the exact WKB connection formula (Paper~4).}
\label{tab:spectral}
\begin{ruledtabular}
\begin{tabular}{l c c c c c c c}
Platform & $\omega/T_H$ & $n_\text{total}$ & $n_\text{Planck}$ & $\delta_\text{disp}$ & $\delta_\text{diss}$ & $\delta_k$ & Dominant \\
\hline
\multicolumn{8}{c}{\textit{Steinhauer $^{87}$Rb}} \\
 & 0.5 & 1.517 & 1.517 & $-8.8\times10^{-5}$ & $1.6\times10^{-7}$ & $3.1\times10^{-7}$ & dispersive \\
 & 1.0 & 0.580 & 0.579 & $-8.8\times10^{-5}$ & $6.1\times10^{-7}$ & $1.2\times10^{-6}$ & dispersive \\
 & 2.0 & 0.157 & 0.157 & $-8.8\times10^{-5}$ & $2.4\times10^{-6}$ & $4.8\times10^{-6}$ & dispersive \\
 & 3.0 & 0.052 & 0.052 & $-8.8\times10^{-5}$ & $5.5\times10^{-6}$ & $1.1\times10^{-5}$ & dispersive \\
 & 5.0 & 0.007 & 0.007 & $-8.8\times10^{-5}$ & $1.5\times10^{-5}$ & $3.0\times10^{-5}$ & dispersive \\
\hline
\multicolumn{8}{c}{\textit{Heidelberg $^{39}$K}} \\
 & 0.5 & 1.517 & 1.517 & $-7.5\times10^{-5}$ & $1.0\times10^{-5}$ & $2.1\times10^{-5}$ & dispersive \\
 & 1.0 & 0.580 & 0.579 & $-7.5\times10^{-5}$ & $4.1\times10^{-5}$ & $8.2\times10^{-5}$ & dispersive \\
 & 2.0 & 0.157 & 0.157 & $-7.5\times10^{-5}$ & $1.6\times10^{-4}$ & $3.2\times10^{-4}$ & noise floor \\
 & 3.0 & 0.052 & 0.052 & $-7.5\times10^{-5}$ & $3.7\times10^{-4}$ & $7.3\times10^{-4}$ & noise floor \\
 & 5.0 & 0.008 & 0.007 & $-7.5\times10^{-5}$ & $1.0\times10^{-3}$ & $2.0\times10^{-3}$ & noise floor \\
\hline
\multicolumn{8}{c}{\textit{Trento $^{23}$Na}} \\
 & 0.5 & 1.517 & 1.517 & $-1.0\times10^{-4}$ & $9.1\times10^{-8}$ & $1.8\times10^{-7}$ & dispersive \\
 & 1.0 & 0.580 & 0.579 & $-1.0\times10^{-4}$ & $3.6\times10^{-7}$ & $7.1\times10^{-7}$ & dispersive \\
 & 2.0 & 0.157 & 0.157 & $-1.0\times10^{-4}$ & $1.4\times10^{-6}$ & $2.8\times10^{-6}$ & dispersive \\
 & 3.0 & 0.052 & 0.052 & $-1.0\times10^{-4}$ & $3.2\times10^{-6}$ & $6.4\times10^{-6}$ & dispersive \\
 & 5.0 & 0.007 & 0.007 & $-1.0\times10^{-4}$ & $8.9\times10^{-6}$ & $1.8\times10^{-5}$ & dispersive \\
\end{tabular}
\end{ruledtabular}
\end{table*}

The key physics features visible in Table~\ref{tab:spectral}:
\begin{itemize}
\item At low frequencies ($\omega \lesssim T_H$), the dispersive correction
dominates. It is negative (subluminal dispersion lowers $T_\text{eff}$).
\item At high frequencies ($\omega \gtrsim 2\,T_H$), the noise floor
dominates for Heidelberg due to its larger damping rate
($\gamma_\text{dim} = 1.6\times10^{-3}$).
\item Steinhauer and Trento have negligible dissipation
($\gamma_\text{dim} \sim 10^{-5}$), so dispersive corrections dominate
across the entire spectrum.
\item The crossover from dispersive to noise-floor dominance occurs at
$\omega/T_H \approx 2$ for Heidelberg only.
\end{itemize}

\section{Detector Requirements}

Three measurement goals are considered:
\begin{enumerate}
\item Detect the first-order dissipative correction $\delta_\text{diss}$
at 3$\sigma$ significance.
\item Detect the FDR noise floor $n_\text{noise}$ at 3$\sigma$.
\item Distinguish exact WKB from perturbative EFT at the 1\% level.
\end{enumerate}

\section{Kappa-Scaling Test}

Varying the surface gravity $\kappa$ while holding BEC material properties
fixed produces different scaling for the two EFT corrections. At fixed
healing length $\xi$, sound speed $c_s$, and transport coefficients
$\gamma_1$, $\gamma_2$:

\begin{align}
|\delta_\text{disp}(\kappa)| &= \frac{\pi}{6}\left(\frac{\xi\kappa}{c_s}\right)^2 \propto \kappa^2 \label{eq:disp_scaling} \\
\delta_\text{diss}(\kappa) &= \frac{(\gamma_1 + \gamma_2)\kappa}{c_s^2} \propto \kappa \label{eq:diss_scaling}
\end{align}

The crossover surface gravity where the two corrections are equal in
magnitude is:
\begin{equation}
\kappa_\text{cross} = \frac{6(\gamma_1 + \gamma_2)}{\pi\xi^2}
\label{eq:crossover}
\end{equation}

This is formally verified in Lean~4 (\texttt{KappaScaling.lean}):
\texttt{kappa\_scaling\_crossover\_balance} proves $|\delta_\text{disp}(\kappa_\text{cross})| = \delta_\text{diss}(\kappa_\text{cross})$;
\texttt{crossover\_unique} proves uniqueness.

\begin{table}[h]
\caption{Kappa-scaling regime classification. $\kappa_\text{nom}$ is the
nominal surface gravity from the transonic background solver.}
\label{tab:kappa_scaling}
\begin{ruledtabular}
\begin{tabular}{lccc}
Platform & $\kappa_\text{nom}/\kappa_\text{cross}$ & Regime & Dominant \\
\hline
Steinhauer & 3.0 & $\kappa > \kappa_\text{cross}$ & dispersive \\
Heidelberg & 0.04 & $\kappa < \kappa_\text{cross}$ & dissipative \\
Trento & 5.7 & $\kappa > \kappa_\text{cross}$ & dispersive \\
\end{tabular}
\end{ruledtabular}
\end{table}

The three platforms span the crossover (Fig.~\ref{fig:kappa_scaling}),
providing complementary experimental access to both scaling regimes.
Heidelberg's large transport coefficients (from the tunable $^{39}$K
scattering length) place it firmly in the dissipative regime, while
Steinhauer and Trento sit in the dispersive regime. Heidelberg is
thus uniquely positioned for the $\kappa$-scaling test: it is the
only platform where the linear (dissipative) scaling is dominant.

\begin{figure}[h]
\includegraphics[width=\columnwidth]{figures/fig46_kappa_scaling_physical.png}
\caption{EFT corrections vs.\ surface gravity for all three platforms.
Dashed lines: $|\delta_\text{disp}| \propto \kappa^2$. Solid lines:
$\delta_\text{diss} \propto \kappa$. Circles: nominal operating point.
Crosses: crossover $\kappa_\text{cross}$ where corrections balance.}
\label{fig:kappa_scaling}
\end{figure}

\section{Polariton Platform Extension (Tier 1)}

Polariton condensates offer an alternative platform with dramatically
higher Hawking temperatures ($T_H \sim 0.06$\,K vs.~$\sim 0.35$\,nK
for BEC) due to the light effective polariton mass
($m^* = 7.0 \times 10^{-35}$\,kg). The Hawking temperature formula
$T_H = \hbar\kappa/(2\pi k_B)$ is unchanged in the driven-dissipative
setting~\cite{Grisins2016}.

The key challenge is the cavity photon decay rate $\Gamma_\text{pol}$,
which introduces frequency-independent spatial attenuation. The
perturbative Tier~1 patch (uniform imaginary frequency shift) is valid
when $\Gamma_\text{pol}/\kappa \ll 1$:

\begin{table}[h]
\caption{Polariton regime classification for $\kappa = 5 \times 10^{10}$\,s$^{-1}$,
$D = 0.1$. Formally verified in \texttt{PolaritonTier1.lean} (6 theorems, zero sorry).}
\label{tab:polariton}
\begin{ruledtabular}
\begin{tabular}{lcccc}
Cavity & $\tau_\text{cav}$ & $\Gamma_\text{pol}/\kappa$ & Regime & EFT test \\
\hline
Standard & 3\,ps & 6.7 & intractable & no \\
Long-lifetime & 100\,ps & 0.20 & borderline & marginal \\
Ultra-long & 300\,ps & 0.067 & perturbative & yes \\
\end{tabular}
\end{ruledtabular}
\end{table}

The spectral degeneracy between polariton decay (flat $\Gamma_\text{pol}$)
and EFT phonon dissipation ($\propto \omega^n$, $n \geq 2$) can be broken
by frequency-resolved measurements: the EFT correction has a distinct
spectral shape that the uniform cavity decay cannot mimic.

Ultra-long-lifetime cavities ($\tau > 300$\,ps, achieved with
Snoke/Pfeiffer samples) place $\Gamma_\text{pol}/\kappa$ in the
perturbative regime, making clean extraction of EFT corrections feasible.
The Paris polariton group (Bramati-Jacquet) has demonstrated horizon
engineering with spatial light modulators, with stimulated Hawking
detection possible 2026--2027.

\section{Summary: Which Platform for Which Measurement}

\begin{table}[h]
\caption{Platform rankings by measurement goal.}
\label{tab:rankings}
\begin{ruledtabular}
\begin{tabular}{lc}
Measurement & Best Platform \\
\hline
Initial EFT detection & Heidelberg \\
FDR noise floor & Heidelberg \\
UV spectral range & Trento \\
Kappa scaling & Heidelberg \\
Dispersive precision & Steinhauer \\
\end{tabular}
\end{ruledtabular}
\end{table}

\end{document}
